\ifdefined\isMainDocument
\else
    \documentclass{article}
    \usepackage{fullpage}
    \usepackage{wrapfig}
    \usepackage{lineno}
    \usepackage{amsmath, amsthm, amsfonts, amssymb, amscd}
    \usepackage{bm, bbm, mathrsfs}
    \usepackage{mathtools}
    \usepackage{nicefrac}
    \usepackage{caption}
    \usepackage{float, multirow}
    \usepackage{textcomp, booktabs}
    \usepackage{biblatex}
    \addbibresource{references.bib}
    \usepackage[colorinlistoftodos]{todonotes}
    \usepackage{graphicx}
    \graphicspath{{../main/}{../supp/}}
    \usepackage{setspace}
    \usepackage{hyperref}
    \setstretch{1.5}
    \begin{document}
    \linenumbers
\fi
% Source: manuscript/scratch.tex (Results). scratch.tex is authoritative; edit
% there first and re-split, or keep the two in step by hand.

\section{Results}
% Numbers moved out of Theory and Model on 2026-10-17 (Theory now carries no data). Place them here
% when Results is built; recheck each against data/processed/ first.
%  - V_A posterior modes 0.003--0.50 (median 0.10): inside the domain of Eq. unlinked_baseline but not
%    uniformly far inside it; a hundredfold deficit would need long-run V_A ~25 at kappa = 1, an
%    extrapolation (label it as the G = 1 extrapolation).
%  - 11 socially monogamous populations, EPP 5--76%, kappa_B = 1.22--1.93; 8 without a pair bond at kappa = 1.
%  - Weak-selection branch: N_e << U_d/(2 V_A) is about 3 for the median population (range 0.6--81);
%    sequence-derived N_e is orders of magnitude larger; 2 N_e s* >= 9,400; Eqs. VA_MSD and HoC agree
%    to better than 1 part in 10^3.
%  - s* median 0.16, range 0.006--0.78 (direct mu), so Q(s) = 2/(1+s) < 2 matters at these values.



\subsection{Intragenomic Variation in Reduction}
While the expected species-wide diversity, $\mathbb{E}[N_e/N]$, provides a global summary of evolutionary constraint, the continuous spatial framework predicts highly heterogeneous topographies of diversity within the genome. By evaluating the Main Equation across independent chromosomes, we can analytically partition the sources of intragenomic variance.

For a focal chromosome $i$, the genome-wide unlinked baseline generated by independent assortment is captured by the denominator of the Main Equation: $1 + Q_u^2(\kappa V_A - \kappa_\text{own}\,v_i/2)$. Because the variance segregating on a single haplotype represents a minor fraction of the genome-wide total or $v_i/2 \ll V_A$, this unlinked constraint is effectively uniform across the genome. Consequently, any predicted variance in standing neutral diversity between chromosomes is driven strictly by the localized, constraint of physical linkage.

Because the localized selective variance is proportional to the functional target size of the chromosome, i.e. $v_i = V_A f_i$, the survival of neutral lineages is governed by the structural ratio of functional density to recombinational map length. Chromosomes with high coding densities and restricted recombination suffer severe exponential decay in local $N_e$, driving rapid lineage sorting. Conversely, chromosomes with expansive map lengths and sparse functional targets effectively isolate neutral alleles from selective backgrounds. These structural dynamics predict the existence of macro-genomic regions that maintain robust neutral diversity even when the global, genome-wide constraint is strong.

\ifdefined\isMainDocument
\else
    \printbibliography
    \end{document}
\fi
